
\chapter{Forecasting and Calibration}
\label{ch:forecasting}

Dispatch safety begins with forecast quality: an optimiser can only be as
reliable as the prediction intervals it conditions on.  This chapter
describes the six model families evaluated in ORIUS, the training
protocol designed to prevent temporal leakage, the conformal prediction
layer that converts point forecasts into calibrated intervals, and the
structural calibration weaknesses that motivate the adaptive inflation
in DC3S.

% ----------------------------------------------------------------
\section{Forecast Targets}

Three operational targets are forecast at hourly resolution over a
48-step rolling horizon ($H = 48$):

\begin{enumerate}
  \item \texttt{load\_mw} --- net system load (MW).
  \item \texttt{wind\_mw} --- aggregate wind generation (MW).
  \item \texttt{solar\_mw} --- aggregate PV generation (MW).
\end{enumerate}

A fourth target, \texttt{price\_eur\_mwh}, is forecast for the Germany
surface only and used in the dispatch objective
(Section~\ref{sec:dispatch-objective}, Chapter~\ref{ch:battery-dynamics}).
All targets are standardised to zero mean and unit variance before model
ingestion; predictions are de-standardised before interval construction.

% ----------------------------------------------------------------
\section{Model Families}
\label{sec:model-families}

Six model families span the spectrum from gradient-boosted trees to
attention-based deep architectures:

\begin{table}[ht]
\centering
\caption{Forecast model families and key hyper-parameters.}
\label{tab:model-families}
\begin{tabular}{llll}
\toprule
\textbf{Family} & \textbf{Type} & \textbf{Key hyper-params} &
  \textbf{Notes} \\
\midrule
GBM        & Gradient-boosted trees & depth, lr, rounds &
  Strongest on locked DE dataset \\
LSTM       & Recurrent (gated)      & hidden, layers, dropout &
  Baseline sequential model \\
TCN        & Temporal convolution   & kernel, dilation, channels &
  Parallelisable alternative to LSTM \\
N-BEATS    & Basis-expansion blocks & stacks, width, sharing &
  Interpretable decomposition \\
TFT        & Temporal fusion        & heads, hidden, quantiles &
  Multi-horizon with variable selection \\
PatchTST   & Patch-based transformer& patch len, heads, depth &
  Channel-independent patching \\
\bottomrule
\end{tabular}
\end{table}

\paragraph{GBM (LightGBM).}
Gradient Boosted Machines trained via LightGBM achieve the lowest MAE on
the locked Germany surface for load and wind targets.  The model ingests
94 engineered features including calendar dummies, rolling statistics, and
weather forecasts.  Quantile regression heads at $\tau \in
\{0.025, 0.05, 0.10, 0.25, 0.50, 0.75, 0.90, 0.95, 0.975\}$ are trained
independently, providing native prediction intervals before conformal
calibration.

\paragraph{LSTM.}
A two-layer LSTM with 128 hidden units and 10\% dropout serves as the
recurrent baseline.  Input sequences of 168\,h (one week) are consumed to
produce a 48-step ahead forecast.  The LSTM captures temporal dependencies
but is prone to over-smoothing during rapid ramp events.

\paragraph{TCN.}
A Temporal Convolutional Network with dilated causal convolutions
(kernel size~7, 4 dilation layers) provides a parallelisable alternative
to the LSTM.  TCN achieves comparable accuracy but trains 3--5$\times$
faster due to the absence of sequential hidden-state propagation.

\paragraph{N-BEATS.}
The Neural Basis Expansion Analysis architecture decomposes the forecast
into trend and seasonality stacks using learnable polynomial and Fourier
basis functions.  The interpretable variant is used, with 3 trend stacks
(polynomial degree~3) and 3 seasonality stacks (harmonics~5).

\paragraph{TFT.}
The Temporal Fusion Transformer combines variable selection, gated residual
connections, and multi-head attention for multi-horizon quantile forecasting.
TFT is the only architecture in the family that natively produces
probabilistic forecasts without a separate quantile-regression head.

\paragraph{PatchTST.}
The Patch Time-Series Transformer segments the input sequence into
non-overlapping patches of length~16, embeds each patch, and applies a
standard transformer encoder with 4 heads and 3 layers.
Channel-independent processing avoids cross-variable leakage and reduces
parameter count.

% ----------------------------------------------------------------
\section{Training Protocol and Leakage Control}
\label{sec:training-protocol}

Temporal leakage is the primary threat to forecast evaluation integrity.
The following protocol prevents it:

\begin{enumerate}
  \item \textbf{Expanding-window split.}  Training data grows monotonically;
    the test window is always strictly future.  No shuffle split is permitted.
  \item \textbf{Embargo gap.}  A 48-hour embargo separates the end of the
    training window from the start of the test window, ensuring that no
    feature derived from a lagged target can leak future information.
  \item \textbf{Walk-forward evaluation.}  The test set is divided into
    non-overlapping 168-hour (one-week) windows.  Metrics are computed per
    window and then aggregated (mean, P95).
  \item \textbf{Feature freeze.}  The feature engineering pipeline is frozen
    before model training begins.  No feature is constructed using information
    from the test period.
  \item \textbf{Seed locking.}  All random seeds (NumPy, PyTorch, LightGBM)
    are fixed at 42 for the primary run, with reproducibility verified across
    seeds $\{11, 22, 33, 44\}$.
\end{enumerate}

% ----------------------------------------------------------------
\section{Conformal Prediction Intervals}
\label{sec:conformal}

Point forecasts are converted to prediction intervals using Conformalized
Quantile Regression (CQR), following the split-conformal framework:

\begin{enumerate}
  \item Train quantile regressors at levels $\alpha/2$ and $1 - \alpha/2$
    on the training fold to obtain $\hat{q}_{\alpha/2}(x)$ and
    $\hat{q}_{1-\alpha/2}(x)$.
  \item On a held-out calibration fold of size~$n$, compute the conformity
    scores:
    \begin{equation}
    \label{eq:cqr-scores}
      s_i = \max\!\bigl(\hat{q}_{\alpha/2}(x_i) - y_i,\;
                        y_i - \hat{q}_{1-\alpha/2}(x_i)\bigr),
      \quad i = 1, \ldots, n.
    \end{equation}
  \item Let $\hat{q}$ be the $\lceil (1-\alpha)(n+1) \rceil / n$
    quantile of $\{s_1, \ldots, s_n\}$.
  \item The prediction interval for a new point~$x_{n+1}$ is:
    \begin{equation}
    \label{eq:cqr-interval}
      C(x_{n+1}) = \bigl[\hat{q}_{\alpha/2}(x_{n+1}) - \hat{q},\;
                         \hat{q}_{1-\alpha/2}(x_{n+1}) + \hat{q}\bigr].
    \end{equation}
\end{enumerate}

CQR provides a finite-sample marginal coverage guarantee:
$\Pr[y_{n+1} \in C(x_{n+1})] \ge 1 - \alpha$, assuming exchangeability
of the calibration and test data.

% ----------------------------------------------------------------
\section{Marginal Coverage}

On the locked Germany surface, the GBM + CQR pipeline achieves
PICP@90 $= 100\%$ and PICP@95 $= 100\%$ for load (marginal coverage).
Wind and solar targets also achieve marginal coverage above the nominal
level.  These aggregate numbers, however, mask conditional failures.

% ----------------------------------------------------------------
\section{Reliability-Conditioned Under-Coverage}

When coverage is stratified by the DC3S reliability score~$w_t$
(Chapter~\ref{ch:dc3s-adapter}), the picture changes substantially.
Table~\ref{tab:cqr-group-coverage} reproduces the locked group-coverage
results.

\begin{table}[ht]
\centering
\caption{CQR group coverage at 90\% nominal level, stratified by
reliability group.  Source:
\texttt{reports/publication/cqr\_group\_coverage.csv}.}
\label{tab:cqr-group-coverage}
\begin{tabular}{llrr}
\toprule
\textbf{Target} & \textbf{Group} & \textbf{PICP@90} &
  \textbf{Mean width (MW)} \\
\midrule
load\_mw  & low  & 0.802 & 2\,147 \\
load\_mw  & mid  & 0.825 & 2\,305 \\
load\_mw  & high & 0.807 & 2\,345 \\
\midrule
wind\_mw  & low  & 0.488 & 464 \\
wind\_mw  & mid  & 0.396 & 285 \\
wind\_mw  & high & 0.410 & 381 \\
\midrule
solar\_mw & low  & 0.903 & 796 \\
solar\_mw & mid  & 0.887 & 869 \\
solar\_mw & high & 0.816 & 853 \\
\bottomrule
\end{tabular}
\end{table}

Wind intervals are catastrophically under-covered in all three groups
(39--49\% vs.\ 90\% nominal), confirming that marginal guarantees do not
transfer to conditional subsets.  Load intervals under-cover by 8--20
percentage points; solar intervals are closest to nominal but still
under-cover in the high-variability group.

% ----------------------------------------------------------------
\section{Solar and Wind Calibration Weaknesses}
\label{sec:calibration-weaknesses}

Two structural factors drive the calibration gap:

\paragraph{Wind non-stationarity.}
Wind generation exhibits regime switches (storm fronts, lulls) that
violate the exchangeability assumption underlying CQR.  The calibration
fold, drawn from a single contiguous period, may not contain sufficient
representation of extreme regimes, leading to narrow intervals that
fail during the test period.

\paragraph{Solar bimodality.}
Solar generation has a deterministic diurnal envelope (sunrise--sunset)
overlaid with a stochastic cloud-cover component.  During nighttime hours,
the true generation is identically zero, making any nonzero interval
width wasteful.  During daytime, cloud transients produce heavy-tailed
residuals that require wide intervals.  A single set of conformity
scores cannot efficiently cover both regimes simultaneously.

% ----------------------------------------------------------------
\section{Why Calibration Alone Cannot Close the Gap}

The group-coverage results demonstrate a fundamental limitation:
\emph{post-hoc conformal calibration on a static calibration fold cannot
guarantee conditional coverage under distribution shift or telemetry
degradation.}

Three pathways to closing the gap are pursued in subsequent chapters:

\begin{enumerate}
  \item \textbf{Reliability-adaptive inflation} (DC3S, Chapter~\ref{ch:dc3s-adapter}):
    the interval width is inflated in real time as a function of the
    observation quality score~$w_t$, the drift flag~$d_t$, and the stale
    count~$s_t$, compensating for the conditional under-coverage.
  \item \textbf{Fault-tolerant invariant tube} (FTIT, Chapter~\ref{ch:dc3s-adapter}):
    a robust state tube propagates worst-case SOC bounds forward, providing
    a safety backstop that does not depend on interval calibration.
  \item \textbf{Online conformal update} (ACI baseline, Chapter~\ref{ch:cpsbench}):
    the conformal quantile is updated online using the most recent
    residuals, partially adapting to distribution shift but without
    observation-quality feedback.
\end{enumerate}

% ----------------------------------------------------------------
\section{Chapter Summary}

This chapter described the six forecast model families (GBM, LSTM, TCN,
N-BEATS, TFT, PatchTST), the leakage-controlled training protocol, and the
CQR conformal calibration layer.  Marginal coverage meets or exceeds the
nominal level on all targets, but reliability-stratified analysis reveals
severe conditional under-coverage, particularly for wind (39--49\% vs.\ 90\%
nominal).  Solar calibration is weakened by diurnal bimodality, and wind by
regime non-stationarity.  These structural calibration failures motivate the
observation-aware adaptive inflation developed in the next chapter.
